\documentclass[12pt]{article}
\usepackage[T1]{fontenc}
\usepackage[utf8]{inputenc}
\usepackage{lmodern}
\usepackage{textcomp}
\usepackage{enumitem}
\usepackage{geometry}
\geometry{margin=1in}
\begin{document}
\begin{center}
Answer Key - Astronomy - Graduate Level Assessment 7 \
\small (Answers are ordered exactly as in the question paper.)
\end{center}

\section*{Multiple Choice}
\begin{enumerate}[label=\arabic*.]
\item \textbf{Answer:} B
\\ \textit{Options:} A) names for the layers of gaseous and metallic hydrogen deep within the planet.; B) alternating bands of rising and falling air at different latitudes.; C) alternating regions of charged particles in Jupiter's magnetic field.; D) regions of the plasma torus created by ions from Io's volcanoes
\\ \textit{Note:} Correct option: B - alternating bands of rising and falling air at different latitudes.
\item \textbf{Answer:} B
\\ \textit{Options:} A) The resonance pulls Io in different directions and generates heat.; B) It makes the orbit of Io slightly elliptical.; C) It creates a gap with no asteriods between the orbits.; D) It prevents formation of the ring material into other moons.
\\ \textit{Note:} Correct option: B - It makes the orbit of Io slightly elliptical.
\item \textbf{Answer:} D
\\ \textit{Options:} A) Tritos and Desmos; B) Tritos and Deimos; C) Phobos and Tritos; D) Phobos and Deimos
\\ \textit{Note:} Correct option: D - Phobos and Deimos
\item \textbf{Answer:} A
\\ \textit{Options:} A) From the dismantling of small moons by impacts; B) From fragments of planets ejected by impacts; C) From dust grains that escape from passing comets; D) From accretion within the solar nebula at the same time the planets formed
\\ \textit{Note:} Correct option: A - From the dismantling of small moons by impacts
\item \textbf{Answer:} A
\\ \textit{Options:} A) Its rotation axis is nearly perpendicular to the plane of the Solar System.; B) It does not have an ozone layer.; C) It does not rotate fast enough.; D) It is too close to the Sun.
\\ \textit{Note:} Correct option: A - Its rotation axis is nearly perpendicular to the plane of the Solar System.
\end{enumerate}

\section*{True / False}
\begin{enumerate}[label=\arabic*.]
\setcounter{enumi}{5}
\item \textbf{Answer:} True
\\ \textit{Options:} A) True; B) False
\\ \textit{Note:} LLM-generated true statement based on MCQ.
\item \textbf{Answer:} True
\\ \textit{Options:} A) True; B) False
\\ \textit{Note:} LLM-generated true statement based on MCQ.
\item \textbf{Answer:} True
\\ \textit{Options:} A) True; B) False
\\ \textit{Note:} LLM-generated true statement based on MCQ.
\item \textbf{Answer:} True
\\ \textit{Options:} A) True; B) False
\\ \textit{Note:} LLM-generated true statement based on MCQ.
\item \textbf{Answer:} True
\\ \textit{Options:} A) True; B) False
\\ \textit{Note:} LLM-generated true statement based on MCQ.
\end{enumerate}

\section*{Long Form Response}
\begin{enumerate}[label=\arabic*.]
\setcounter{enumi}{10}
\item \textbf{Answer:} Hawking radiation. Explain why this is the correct answer and provide supporting evidence. Reference key facts or concepts that support this choice.
\\ \textit{Note:} Long-form derived from MCQ; award credit for stating the correct idea and giving supporting reasoning.
\item \textbf{Answer:} Orion. Explain why this is the correct answer and provide supporting evidence. Reference key facts or concepts that support this choice.
\\ \textit{Note:} Long-form derived from MCQ; award credit for stating the correct idea and giving supporting reasoning.
\end{enumerate}

\vfill
\noindent\textit{End of Answer Key}
\end{document}